% answers_short.tex —— 2003 年·全国新课程卷·理科 · 简短答案速查
% 由 tools/make_short_answers.py 自 answers.tex 的 \gksolhead/\gkanswer 提取，
% 并与 2003/2003_数学真题_全国新课程理_答案要点.md 逐题比对后生成，勿手改（改源 MD 后重跑）。

\begin{gknotice}
\textbf{使用说明}：本节只给最终答案便于快速核对，完整推导见解析版。答案为编写者独立推导，\textbf{非官方标准答案}。
\end{gknotice}

\gkpart{选择题}
\begin{center}\begin{tabular}{ccccccccccccc}
\toprule
题号 & 1 & 2 & 3 & 4 & 5 & 6 & 7 & 8 & 9 & 10 & 11 & 12 \\
\midrule
答案 & B & D & D & B & B & C & B & C & D & C & B & A \\
\bottomrule
\end{tabular}\end{center}

\gkpart{填空题}
\begin{itemize}[leftmargin=2.6em]
  \item[\textbf{13.}] $-\dfrac{21}{2}$
  \item[\textbf{14.}] $6,30,10$
  \item[\textbf{15.}] $120$
  \item[\textbf{16.}] ②⑤
\end{itemize}

\gkpart{解答题}
\begin{itemize}[leftmargin=2.6em]
  \item[\textbf{17.}] (1) 最小正周期为 $\pi$，最大值为 $1+\sqrt2$；(2) 图象见解析版
  \item[\textbf{18.}] (1) $\arcsin\dfrac{\sqrt2}{3}$；(2) $\dfrac{2\sqrt6}{3}$
  \item[\textbf{19.}] 在 $0<a<1$ 时按 $(1\pm\sqrt{1-a})^2$ 分段单调；$a\geqslant1$ 时在 $(0,+\infty)$ 上单调递增。
  \item[\textbf{20.}] (1) $\xi$ 的分布为 $\left(\dfrac3{25},\dfrac25,\dfrac{28}{75},\dfrac8{75}\right)$；$\eta=3-\xi$；(2) $E\xi=\dfrac{22}{15},\ E\eta=\dfrac{23}{15}$
  \item[\textbf{21.}] 见解析版；轨迹为椭圆，焦点分类按 $a$ 与 $1/\sqrt2$ 的大小确定。
  \item[\textbf{22.}] (1) $a_2=4,\ a_3=13$；$a_n=\dfrac15[3^n+(-1)^{n-1}2^n]+(-1)^n2^na_0$；(2) $0<a_0<\dfrac13$
\end{itemize}
